% !TEX program = xelatex
\documentclass[a4paper,11pt,openany,fontset=none]{ctexrep}

\newcommand{\reporttitle}{5\%总回报目标：开放研究宇宙框架}
\newcommand{\reportsubtitle}{规模中性研究、逐项尽调与动态组合审批}
\newcommand{\reportkicker}{A股多资产总回报目标研究}
\newcommand{\reportscope}{开放研究宇宙 | 直接股票 | ETF | 公募REIT | 债券与现金工具 | 研究用途}
\newcommand{\reportdate}{2026年8月14日}
\newcommand{\reportdatacutoff}{公开规则、产品报告及分配公告截至2026年8月12日；实时价格、流动性、税费与账户资格未收集}
\newcommand{\reporttype}{战略配置研究框架（实施未获批准）}
\newcommand{\reportauthor}{AStock Research}
\newcommand{\reporthouseview}{研究框架采用比例口径、与资金规模无关，不预设任何账户层事实。研究宇宙向所有合规可及的直接股票、ETF、公募REIT、债券与现金工具开放；现有十只产品只是首轮尽调队列，并非排他菜单。唯一的策略研究目标是滚动五年、成本后年化TWRR达到5\%；它不是每年最低收益或收益保证。现金分配只服务于总回报核算和产品尽调，不构成独立目标。当前批准可投池和批准组合均为空。}
\newcommand{\reportquality}{本报告使用证监会、上交所和基金管理人的原始规则、年报与分配公告，支持首轮产品的历史或已公告事实及产品机制。它们不支持未来收益率、年度下限、策略命中率、概率或实施资格；任何非零持仓均须补齐逐项权益研究级尽调、组合风险预算、适当性和实时资料。}
\newcommand{\reportdisclaimer}{本报告基于公开资料整理，仅作研究与教育参考，不构成证券买卖建议，不保证收益或资产价值安全。}

% Reusable institutional-style preamble for XeLaTeX AStock research reports.
% Usage: define \reporttitle and \reportdate, then % Reusable institutional-style preamble for XeLaTeX AStock research reports.
% Usage: define \reporttitle and \reportdate, then % Reusable institutional-style preamble for XeLaTeX AStock research reports.
% Usage: define \reporttitle and \reportdate, then \input{.agents/templates/preamble.tex}.
% Design benchmark: JPM first-page dashboard, BlackRock house-view framing,
% Vanguard scenario discipline, and AB/JPM chartbook exhibit structure.

\providecommand{\reporttitle}{AStock Research Report}
\providecommand{\reportsubtitle}{House View, Evidence, Valuation and Risk}
\providecommand{\reportkicker}{INSTITUTIONAL EQUITY RESEARCH}
\providecommand{\reportscope}{CHINA A-SHARES}
\providecommand{\reportdate}{\today}
\providecommand{\reportdatacutoff}{Data cutoff not specified}
\providecommand{\reporttype}{Research note}
\providecommand{\reportauthor}{AStock Research Agent}
\providecommand{\reporthouseview}{State the house view in one decision-useful paragraph.}
\providecommand{\reportquality}{Source quality not specified.}
\providecommand{\reportdisclaimer}{本报告基于公开资料整理，不构成任何证券买卖建议。}

% ============ 页面布局 ============
\usepackage[top=1.8cm, bottom=1.9cm, left=1.8cm, right=1.8cm]{geometry}
\setlength{\parskip}{0.3em}
\setlength{\parindent}{2em}
\renewcommand{\baselinestretch}{1.12}
\raggedbottom

% ============ Fonts ============
% Prefer the installed Noto CJK family when available (for Linux builds).
% Otherwise retain the original MacTeX typography. This avoids platform
% probing via shell escape and keeps the source portable across both engines.
\IfFontExistsTF{Noto Serif CJK SC}{
  \setCJKmainfont{Noto Serif CJK SC}
  \setCJKsansfont{Noto Sans CJK SC}
  \setCJKmonofont{Noto Sans CJK SC}
  \setmainfont{Liberation Serif}
  \newcommand{\kaishu}{\CJKfontspec{Noto Serif CJK SC}}
  \newcommand{\heiti}{\CJKfontspec{Noto Sans CJK SC}}
  \newcommand{\songti}{\CJKfontspec{Noto Serif CJK SC}}
}{
  \setCJKmainfont[BoldFont={Heiti SC}]{STSong}
  \setCJKsansfont{Heiti SC}
  \setCJKmonofont{Heiti SC}
  \setmainfont{Times New Roman}
  \newcommand{\kaishu}{\CJKfontspec{STKaiti}}
  \newcommand{\heiti}{\CJKfontspec{Heiti SC}}
  \newcommand{\songti}{\CJKfontspec{STSong}}
}

% ============ 颜色（IB Navy/Grey palette） ============
\usepackage{xcolor}
\definecolor{navy}{HTML}{003366}
\definecolor{deepnavy}{HTML}{0B1F33}
\definecolor{steelgrey}{HTML}{4A5568}
\definecolor{mutedgrey}{HTML}{718096}
\definecolor{lightgrey}{HTML}{F7FAFC}
\definecolor{panelgrey}{HTML}{EDF2F7}
\definecolor{riskred}{HTML}{C53030}
\definecolor{riskamber}{HTML}{D69E2E}
\definecolor{riskgreen}{HTML}{38A169}
\definecolor{riskblue}{HTML}{2B6CB0}
\definecolor{accentblue}{HTML}{2B6CB0}
\definecolor{nvgreen}{HTML}{76B900}
\definecolor{rulegrey}{HTML}{CBD5E0}

% ============ 表格 ============
\usepackage{booktabs}
\usepackage{longtable}
\usepackage{tabularx}
\usepackage{multirow}
\usepackage{array}
\usepackage{makecell}
\usepackage{colortbl}
\renewcommand{\arraystretch}{1.12}
\newcolumntype{Y}{>{\centering\arraybackslash}X}
\newcolumntype{L}[1]{>{\raggedright\arraybackslash}p{#1}}
\newcolumntype{C}[1]{>{\centering\arraybackslash}p{#1}}
\newcolumntype{R}[1]{>{\raggedleft\arraybackslash}p{#1}}

% ============ 图形 ============
\usepackage{graphicx}
\usepackage{tikz}
\usetikzlibrary{shapes,arrows.meta,positioning,calc,backgrounds,fit,decorations.pathreplacing}
\usepackage{pgfplots}
\usepgfplotslibrary{polar}
\pgfplotsset{compat=1.18}
\usepackage[edges]{forest}

% ============ 页眉页脚（报告标题和日期需在main.tex中覆盖） ============
\usepackage{fancyhdr}
\pagestyle{fancy}
\fancyhf{}
\fancyhead[L]{\scriptsize\color{steelgrey} \reporttitle}
\fancyhead[R]{\scriptsize\color{steelgrey} \reportdate}
\fancyfoot[C]{\small\color{steelgrey} \thepage}
\fancyfoot[R]{\tiny\color{steelgrey} \reportdisclaimer}
\renewcommand{\headrulewidth}{0.4pt}
\renewcommand{\footrulewidth}{0.2pt}

% ============ 章节格式 ============
\usepackage{titlesec}
\titleformat{\chapter}[block]
  {\normalfont\LARGE\bfseries\color{navy}}
  {\thechapter}{0.9em}{}
\titlespacing*{\chapter}{0pt}{0.4em}{1.0em}
\titleformat{\section}
  {\normalfont\large\bfseries\color{navy}}
  {\thesection}{1em}{}
\titleformat{\subsection}
  {\normalfont\normalsize\bfseries\color{steelgrey}}
  {\thesubsection}{1em}{}

% ============ 超链接 ============
\usepackage{hyperref}
\hypersetup{
  colorlinks=true,
  linkcolor=navy,
  urlcolor=navy,
  citecolor=navy,
}

% ============ 其他工具 ============
\usepackage{enumitem}
\setlist{nosep, leftmargin=2em}
\IfFileExists{fontawesome5.sty}{\usepackage{fontawesome5}}{}
% ============ Class B 全局兜底（xeCJK + 缺字 fallback 阻断）：
% 当 fontawesome5 因 Times New Roman 缺字 / \FA 命名空间失败时，
% 禁止输出 aExclamationTriangle / aCloud 裸文本，一律降级为 amssymb 原生符号。
% 注：amssymb 在 L112 已加载，此处不重复。所有 \def 在本段落。
\def\faExclamationTriangle{\ensuremath{\textcolor{riskred}{\blacktriangle}}\kern-0.25em}
\def\faCloud{\ensuremath{\clubsuit}\kern-0.1em}
\def\faYenSign{\ensuremath{\yen}}
\def\faIndustry{\ensuremath{\S}}
\def\faCogs{\ensuremath{\ast}}
\def\faFlagCheckered{\ensuremath{\blacktriangleright}}
\def\faMemory{\textbf{[RAM]}}
\def\faServer{\textbf{[DC]}}
\def\faGlobe{\ensuremath{\diamond}}
\def\faStar{\ensuremath{\bigstar}}
\def\faMicrochip{\textbf{[$\mu$P]}}
\def\faLayerGroup{\textbf{[LAY]}}
\def\faCubes{\ensuremath{\boxtimes}}
\usepackage{amssymb}
\usepackage{pifont}
\usepackage{tcolorbox}
\tcbuselibrary{skins,breakable}
\usepackage{caption}
\captionsetup{font=small, skip=4pt}
\usepackage{float}
\usepackage{needspace}
\setlength{\textfloatsep}{8pt plus 2pt minus 2pt}
\setlength{\floatsep}{6pt plus 2pt minus 2pt}
\setlength{\intextsep}{6pt plus 2pt minus 2pt}
\setlength{\abovecaptionskip}{4pt}
\setlength{\belowcaptionskip}{2pt}

% ============ 自定义命令 ============
\newcommand{\rmark}{\textcolor{riskred}{\ding{108}}}
\newcommand{\amark}{\textcolor{riskamber}{\ding{108}}}
\newcommand{\gmark}{\textcolor{riskgreen}{\ding{108}}}
\newcommand{\starfive}{$\bigstar\bigstar\bigstar\bigstar\bigstar$}
\newcommand{\starfour}{$\bigstar\bigstar\bigstar\bigstar$}
\newcommand{\starthree}{$\bigstar\bigstar\bigstar$}
% 资料来源注：使用楷体（STKaiti）而非斜体，避免中文无斜体字形导致的字体回退警告
% \mdseries 避免继承周围 \bfseries 导致 STKaiti 粗体缺失警告
% \sloppy + \hbadness 允许长文件名/URL 在必要时断行
% 局部重定义 \textbf → \heiti：STKaiti 无粗体，用黑体替代强调
\newcommand{\sourcenote}[1]{\par{\scriptsize\mdseries\color{mutedgrey}\kaishu\sloppy\hbadness=9999\def\textbf##1{{\heiti ##1}} 资料来源：#1}\vspace{0.2em}}
\newcommand{\qualitynote}[1]{\par{\scriptsize\mdseries\color{mutedgrey}\sloppy\def\textbf##1{{\heiti ##1}} 数据质量：#1}\vspace{0.2em}}
\newcommand{\evidencetag}[1]{%
  \tikz[baseline=(x.base)]\node[
    xshift=0pt,
    inner xsep=3.5pt,
    inner ysep=1.2pt,
    rounded corners=1.5pt,
    fill=panelgrey,
    draw=rulegrey,
    text=steelgrey,
    font=\scriptsize
  ] (x) {#1};%
}
\newcommand{\stance}[2]{\textbf{\textcolor{#1}{#2}}}

% 信息框
\newtcolorbox{keyinsight}[1][]{
  colback=lightgrey, colframe=navy,
  fonttitle=\bfseries, title={#1}, breakable,
  arc=2pt, boxrule=0.5pt, left=8pt, right=8pt, top=6pt, bottom=6pt
}
\newtcolorbox{riskbox}[1][]{
  colback=red!5, colframe=riskred,
  fonttitle=\bfseries, title={#1}, breakable,
  arc=2pt, boxrule=0.5pt, left=8pt, right=8pt, top=6pt, bottom=6pt
}
\newtcolorbox{dashboardbox}[1][]{
  enhanced, breakable,
  colback=white, colframe=navy,
  boxrule=0.7pt, arc=2pt,
  fonttitle=\bfseries\color{white},
  coltitle=white, colbacktitle=navy,
  title={#1},
  left=9pt, right=9pt, top=7pt, bottom=7pt
}
\newtcolorbox{houseviewbox}[1][]{
  enhanced, breakable,
  colback=lightgrey, colframe=deepnavy,
  borderline west={3pt}{0pt}{accentblue},
  boxrule=0.4pt, arc=2pt,
  fonttitle=\bfseries, title={#1},
  left=9pt, right=9pt, top=7pt, bottom=7pt
}
\newtcolorbox{sourcequalitybox}[1][]{
  enhanced, breakable,
  colback=yellow!6, colframe=riskamber,
  boxrule=0.45pt, arc=2pt,
  fonttitle=\bfseries, title={#1},
  left=9pt, right=9pt, top=7pt, bottom=7pt
}
% exhibitbox：表/图外框。left=7pt+right=7pt=14pt 内边距，before upper 内的 \linewidth 已扣除该内边距。
% 约定：\exhibitboxwidth = 真实内边距宽（tcolorbox 已扣好 padding 的 \linewidth）。
% 由于 tabularx/longtable 宽度参数的语义是「外部总宽（含两侧 \tabcolsep）」，内部列实际可用 = W - 2\tabcolsep。
% 所以 tabularx 直接传 \exhibitboxwidth（总宽 == 内边距宽），而其他用途（resizebox、tikz、p{...}）用 \exhibitwidth = \exhibitboxwidth - 2\tabcolsep。
\newlength{\exhibitboxwidth}
\newlength{\exhibitwidth}
\newtcolorbox{exhibitbox}[1][]{
  enhanced, breakable,
  colback=white, colframe=rulegrey,
  borderline north={1.2pt}{0pt}{navy},
  boxrule=0.35pt, arc=1pt,
  fonttitle=\bfseries\color{navy}, title={#1},
  left=7pt, right=7pt, top=6pt, bottom=6pt,
  before upper={\global\exhibitboxwidth=\linewidth\relax
    \global\exhibitwidth=\dimexpr\linewidth-2\tabcolsep\relax},
}
\newtcolorbox{disclosurebox}[1][]{
  enhanced, breakable,
  colback=lightgrey, colframe=steelgrey,
  boxrule=0.35pt, arc=1pt,
  fonttitle=\bfseries, title={#1},
  left=8pt, right=8pt, top=6pt, bottom=6pt
}

% ============ First-page dashboard cover ============
\newcommand{\astockcover}{%
\begin{titlepage}
\thispagestyle{empty}
\begin{tikzpicture}[remember picture, overlay]
  \fill[deepnavy] (current page.north west) rectangle ([yshift=-2.55cm]current page.north east);
  \fill[accentblue] ([yshift=-2.55cm]current page.north west) rectangle ([yshift=-2.66cm]current page.north east);
  \fill[deepnavy] (current page.south west) rectangle ([yshift=1.2cm]current page.south east);
\end{tikzpicture}
\vspace*{0.62cm}
{\color{white}\sffamily\bfseries\small \reportkicker}\\[0.18cm]
{\color{steelgrey}\sffamily\scriptsize \reportscope \quad|\quad \reporttype \quad|\quad \reportdate}\\[1.0cm]
{\fontsize{25}{30}\selectfont\bfseries\color{navy}\reporttitle}\\[0.25cm]
{\fontsize{13}{17}\selectfont\color{steelgrey}\reportsubtitle}\\[0.75cm]
\begin{dashboardbox}[决策摘要]
\begin{tabularx}{\textwidth}{L{3.0cm}X L{3.0cm}X}
\textbf{报告类型} & \reporttype & \textbf{数据截至} & \reportdatacutoff \\
\textbf{覆盖范围} & \reportscope & \textbf{作者} & \reportauthor \\
\end{tabularx}
\vspace{0.45em}
\begin{houseviewbox}[本机构观点]
\reporthouseview
\end{houseviewbox}
\begin{sourcequalitybox}[证据质量]
\reportquality
\end{sourcequalitybox}
\end{dashboardbox}
\vfill
{\scriptsize\color{steelgrey}\reportdisclaimer}
\end{titlepage}
}
.
% Design benchmark: JPM first-page dashboard, BlackRock house-view framing,
% Vanguard scenario discipline, and AB/JPM chartbook exhibit structure.

\providecommand{\reporttitle}{AStock Research Report}
\providecommand{\reportsubtitle}{House View, Evidence, Valuation and Risk}
\providecommand{\reportkicker}{INSTITUTIONAL EQUITY RESEARCH}
\providecommand{\reportscope}{CHINA A-SHARES}
\providecommand{\reportdate}{\today}
\providecommand{\reportdatacutoff}{Data cutoff not specified}
\providecommand{\reporttype}{Research note}
\providecommand{\reportauthor}{AStock Research Agent}
\providecommand{\reporthouseview}{State the house view in one decision-useful paragraph.}
\providecommand{\reportquality}{Source quality not specified.}
\providecommand{\reportdisclaimer}{本报告基于公开资料整理，不构成任何证券买卖建议。}

% ============ 页面布局 ============
\usepackage[top=1.8cm, bottom=1.9cm, left=1.8cm, right=1.8cm]{geometry}
\setlength{\parskip}{0.3em}
\setlength{\parindent}{2em}
\renewcommand{\baselinestretch}{1.12}
\raggedbottom

% ============ Fonts ============
% Prefer the installed Noto CJK family when available (for Linux builds).
% Otherwise retain the original MacTeX typography. This avoids platform
% probing via shell escape and keeps the source portable across both engines.
\IfFontExistsTF{Noto Serif CJK SC}{
  \setCJKmainfont{Noto Serif CJK SC}
  \setCJKsansfont{Noto Sans CJK SC}
  \setCJKmonofont{Noto Sans CJK SC}
  \setmainfont{Liberation Serif}
  \newcommand{\kaishu}{\CJKfontspec{Noto Serif CJK SC}}
  \newcommand{\heiti}{\CJKfontspec{Noto Sans CJK SC}}
  \newcommand{\songti}{\CJKfontspec{Noto Serif CJK SC}}
}{
  \setCJKmainfont[BoldFont={Heiti SC}]{STSong}
  \setCJKsansfont{Heiti SC}
  \setCJKmonofont{Heiti SC}
  \setmainfont{Times New Roman}
  \newcommand{\kaishu}{\CJKfontspec{STKaiti}}
  \newcommand{\heiti}{\CJKfontspec{Heiti SC}}
  \newcommand{\songti}{\CJKfontspec{STSong}}
}

% ============ 颜色（IB Navy/Grey palette） ============
\usepackage{xcolor}
\definecolor{navy}{HTML}{003366}
\definecolor{deepnavy}{HTML}{0B1F33}
\definecolor{steelgrey}{HTML}{4A5568}
\definecolor{mutedgrey}{HTML}{718096}
\definecolor{lightgrey}{HTML}{F7FAFC}
\definecolor{panelgrey}{HTML}{EDF2F7}
\definecolor{riskred}{HTML}{C53030}
\definecolor{riskamber}{HTML}{D69E2E}
\definecolor{riskgreen}{HTML}{38A169}
\definecolor{riskblue}{HTML}{2B6CB0}
\definecolor{accentblue}{HTML}{2B6CB0}
\definecolor{nvgreen}{HTML}{76B900}
\definecolor{rulegrey}{HTML}{CBD5E0}

% ============ 表格 ============
\usepackage{booktabs}
\usepackage{longtable}
\usepackage{tabularx}
\usepackage{multirow}
\usepackage{array}
\usepackage{makecell}
\usepackage{colortbl}
\renewcommand{\arraystretch}{1.12}
\newcolumntype{Y}{>{\centering\arraybackslash}X}
\newcolumntype{L}[1]{>{\raggedright\arraybackslash}p{#1}}
\newcolumntype{C}[1]{>{\centering\arraybackslash}p{#1}}
\newcolumntype{R}[1]{>{\raggedleft\arraybackslash}p{#1}}

% ============ 图形 ============
\usepackage{graphicx}
\usepackage{tikz}
\usetikzlibrary{shapes,arrows.meta,positioning,calc,backgrounds,fit,decorations.pathreplacing}
\usepackage{pgfplots}
\usepgfplotslibrary{polar}
\pgfplotsset{compat=1.18}
\usepackage[edges]{forest}

% ============ 页眉页脚（报告标题和日期需在main.tex中覆盖） ============
\usepackage{fancyhdr}
\pagestyle{fancy}
\fancyhf{}
\fancyhead[L]{\scriptsize\color{steelgrey} \reporttitle}
\fancyhead[R]{\scriptsize\color{steelgrey} \reportdate}
\fancyfoot[C]{\small\color{steelgrey} \thepage}
\fancyfoot[R]{\tiny\color{steelgrey} \reportdisclaimer}
\renewcommand{\headrulewidth}{0.4pt}
\renewcommand{\footrulewidth}{0.2pt}

% ============ 章节格式 ============
\usepackage{titlesec}
\titleformat{\chapter}[block]
  {\normalfont\LARGE\bfseries\color{navy}}
  {\thechapter}{0.9em}{}
\titlespacing*{\chapter}{0pt}{0.4em}{1.0em}
\titleformat{\section}
  {\normalfont\large\bfseries\color{navy}}
  {\thesection}{1em}{}
\titleformat{\subsection}
  {\normalfont\normalsize\bfseries\color{steelgrey}}
  {\thesubsection}{1em}{}

% ============ 超链接 ============
\usepackage{hyperref}
\hypersetup{
  colorlinks=true,
  linkcolor=navy,
  urlcolor=navy,
  citecolor=navy,
}

% ============ 其他工具 ============
\usepackage{enumitem}
\setlist{nosep, leftmargin=2em}
\IfFileExists{fontawesome5.sty}{\usepackage{fontawesome5}}{}
% ============ Class B 全局兜底（xeCJK + 缺字 fallback 阻断）：
% 当 fontawesome5 因 Times New Roman 缺字 / \FA 命名空间失败时，
% 禁止输出 aExclamationTriangle / aCloud 裸文本，一律降级为 amssymb 原生符号。
% 注：amssymb 在 L112 已加载，此处不重复。所有 \def 在本段落。
\def\faExclamationTriangle{\ensuremath{\textcolor{riskred}{\blacktriangle}}\kern-0.25em}
\def\faCloud{\ensuremath{\clubsuit}\kern-0.1em}
\def\faYenSign{\ensuremath{\yen}}
\def\faIndustry{\ensuremath{\S}}
\def\faCogs{\ensuremath{\ast}}
\def\faFlagCheckered{\ensuremath{\blacktriangleright}}
\def\faMemory{\textbf{[RAM]}}
\def\faServer{\textbf{[DC]}}
\def\faGlobe{\ensuremath{\diamond}}
\def\faStar{\ensuremath{\bigstar}}
\def\faMicrochip{\textbf{[$\mu$P]}}
\def\faLayerGroup{\textbf{[LAY]}}
\def\faCubes{\ensuremath{\boxtimes}}
\usepackage{amssymb}
\usepackage{pifont}
\usepackage{tcolorbox}
\tcbuselibrary{skins,breakable}
\usepackage{caption}
\captionsetup{font=small, skip=4pt}
\usepackage{float}
\usepackage{needspace}
\setlength{\textfloatsep}{8pt plus 2pt minus 2pt}
\setlength{\floatsep}{6pt plus 2pt minus 2pt}
\setlength{\intextsep}{6pt plus 2pt minus 2pt}
\setlength{\abovecaptionskip}{4pt}
\setlength{\belowcaptionskip}{2pt}

% ============ 自定义命令 ============
\newcommand{\rmark}{\textcolor{riskred}{\ding{108}}}
\newcommand{\amark}{\textcolor{riskamber}{\ding{108}}}
\newcommand{\gmark}{\textcolor{riskgreen}{\ding{108}}}
\newcommand{\starfive}{$\bigstar\bigstar\bigstar\bigstar\bigstar$}
\newcommand{\starfour}{$\bigstar\bigstar\bigstar\bigstar$}
\newcommand{\starthree}{$\bigstar\bigstar\bigstar$}
% 资料来源注：使用楷体（STKaiti）而非斜体，避免中文无斜体字形导致的字体回退警告
% \mdseries 避免继承周围 \bfseries 导致 STKaiti 粗体缺失警告
% \sloppy + \hbadness 允许长文件名/URL 在必要时断行
% 局部重定义 \textbf → \heiti：STKaiti 无粗体，用黑体替代强调
\newcommand{\sourcenote}[1]{\par{\scriptsize\mdseries\color{mutedgrey}\kaishu\sloppy\hbadness=9999\def\textbf##1{{\heiti ##1}} 资料来源：#1}\vspace{0.2em}}
\newcommand{\qualitynote}[1]{\par{\scriptsize\mdseries\color{mutedgrey}\sloppy\def\textbf##1{{\heiti ##1}} 数据质量：#1}\vspace{0.2em}}
\newcommand{\evidencetag}[1]{%
  \tikz[baseline=(x.base)]\node[
    xshift=0pt,
    inner xsep=3.5pt,
    inner ysep=1.2pt,
    rounded corners=1.5pt,
    fill=panelgrey,
    draw=rulegrey,
    text=steelgrey,
    font=\scriptsize
  ] (x) {#1};%
}
\newcommand{\stance}[2]{\textbf{\textcolor{#1}{#2}}}

% 信息框
\newtcolorbox{keyinsight}[1][]{
  colback=lightgrey, colframe=navy,
  fonttitle=\bfseries, title={#1}, breakable,
  arc=2pt, boxrule=0.5pt, left=8pt, right=8pt, top=6pt, bottom=6pt
}
\newtcolorbox{riskbox}[1][]{
  colback=red!5, colframe=riskred,
  fonttitle=\bfseries, title={#1}, breakable,
  arc=2pt, boxrule=0.5pt, left=8pt, right=8pt, top=6pt, bottom=6pt
}
\newtcolorbox{dashboardbox}[1][]{
  enhanced, breakable,
  colback=white, colframe=navy,
  boxrule=0.7pt, arc=2pt,
  fonttitle=\bfseries\color{white},
  coltitle=white, colbacktitle=navy,
  title={#1},
  left=9pt, right=9pt, top=7pt, bottom=7pt
}
\newtcolorbox{houseviewbox}[1][]{
  enhanced, breakable,
  colback=lightgrey, colframe=deepnavy,
  borderline west={3pt}{0pt}{accentblue},
  boxrule=0.4pt, arc=2pt,
  fonttitle=\bfseries, title={#1},
  left=9pt, right=9pt, top=7pt, bottom=7pt
}
\newtcolorbox{sourcequalitybox}[1][]{
  enhanced, breakable,
  colback=yellow!6, colframe=riskamber,
  boxrule=0.45pt, arc=2pt,
  fonttitle=\bfseries, title={#1},
  left=9pt, right=9pt, top=7pt, bottom=7pt
}
% exhibitbox：表/图外框。left=7pt+right=7pt=14pt 内边距，before upper 内的 \linewidth 已扣除该内边距。
% 约定：\exhibitboxwidth = 真实内边距宽（tcolorbox 已扣好 padding 的 \linewidth）。
% 由于 tabularx/longtable 宽度参数的语义是「外部总宽（含两侧 \tabcolsep）」，内部列实际可用 = W - 2\tabcolsep。
% 所以 tabularx 直接传 \exhibitboxwidth（总宽 == 内边距宽），而其他用途（resizebox、tikz、p{...}）用 \exhibitwidth = \exhibitboxwidth - 2\tabcolsep。
\newlength{\exhibitboxwidth}
\newlength{\exhibitwidth}
\newtcolorbox{exhibitbox}[1][]{
  enhanced, breakable,
  colback=white, colframe=rulegrey,
  borderline north={1.2pt}{0pt}{navy},
  boxrule=0.35pt, arc=1pt,
  fonttitle=\bfseries\color{navy}, title={#1},
  left=7pt, right=7pt, top=6pt, bottom=6pt,
  before upper={\global\exhibitboxwidth=\linewidth\relax
    \global\exhibitwidth=\dimexpr\linewidth-2\tabcolsep\relax},
}
\newtcolorbox{disclosurebox}[1][]{
  enhanced, breakable,
  colback=lightgrey, colframe=steelgrey,
  boxrule=0.35pt, arc=1pt,
  fonttitle=\bfseries, title={#1},
  left=8pt, right=8pt, top=6pt, bottom=6pt
}

% ============ First-page dashboard cover ============
\newcommand{\astockcover}{%
\begin{titlepage}
\thispagestyle{empty}
\begin{tikzpicture}[remember picture, overlay]
  \fill[deepnavy] (current page.north west) rectangle ([yshift=-2.55cm]current page.north east);
  \fill[accentblue] ([yshift=-2.55cm]current page.north west) rectangle ([yshift=-2.66cm]current page.north east);
  \fill[deepnavy] (current page.south west) rectangle ([yshift=1.2cm]current page.south east);
\end{tikzpicture}
\vspace*{0.62cm}
{\color{white}\sffamily\bfseries\small \reportkicker}\\[0.18cm]
{\color{steelgrey}\sffamily\scriptsize \reportscope \quad|\quad \reporttype \quad|\quad \reportdate}\\[1.0cm]
{\fontsize{25}{30}\selectfont\bfseries\color{navy}\reporttitle}\\[0.25cm]
{\fontsize{13}{17}\selectfont\color{steelgrey}\reportsubtitle}\\[0.75cm]
\begin{dashboardbox}[决策摘要]
\begin{tabularx}{\textwidth}{L{3.0cm}X L{3.0cm}X}
\textbf{报告类型} & \reporttype & \textbf{数据截至} & \reportdatacutoff \\
\textbf{覆盖范围} & \reportscope & \textbf{作者} & \reportauthor \\
\end{tabularx}
\vspace{0.45em}
\begin{houseviewbox}[本机构观点]
\reporthouseview
\end{houseviewbox}
\begin{sourcequalitybox}[证据质量]
\reportquality
\end{sourcequalitybox}
\end{dashboardbox}
\vfill
{\scriptsize\color{steelgrey}\reportdisclaimer}
\end{titlepage}
}
.
% Design benchmark: JPM first-page dashboard, BlackRock house-view framing,
% Vanguard scenario discipline, and AB/JPM chartbook exhibit structure.

\providecommand{\reporttitle}{AStock Research Report}
\providecommand{\reportsubtitle}{House View, Evidence, Valuation and Risk}
\providecommand{\reportkicker}{INSTITUTIONAL EQUITY RESEARCH}
\providecommand{\reportscope}{CHINA A-SHARES}
\providecommand{\reportdate}{\today}
\providecommand{\reportdatacutoff}{Data cutoff not specified}
\providecommand{\reporttype}{Research note}
\providecommand{\reportauthor}{AStock Research Agent}
\providecommand{\reporthouseview}{State the house view in one decision-useful paragraph.}
\providecommand{\reportquality}{Source quality not specified.}
\providecommand{\reportdisclaimer}{本报告基于公开资料整理，不构成任何证券买卖建议。}

% ============ 页面布局 ============
\usepackage[top=1.8cm, bottom=1.9cm, left=1.8cm, right=1.8cm]{geometry}
\setlength{\parskip}{0.3em}
\setlength{\parindent}{2em}
\renewcommand{\baselinestretch}{1.12}
\raggedbottom

% ============ Fonts ============
% Prefer the installed Noto CJK family when available (for Linux builds).
% Otherwise retain the original MacTeX typography. This avoids platform
% probing via shell escape and keeps the source portable across both engines.
\IfFontExistsTF{Noto Serif CJK SC}{
  \setCJKmainfont{Noto Serif CJK SC}
  \setCJKsansfont{Noto Sans CJK SC}
  \setCJKmonofont{Noto Sans CJK SC}
  \setmainfont{Liberation Serif}
  \newcommand{\kaishu}{\CJKfontspec{Noto Serif CJK SC}}
  \newcommand{\heiti}{\CJKfontspec{Noto Sans CJK SC}}
  \newcommand{\songti}{\CJKfontspec{Noto Serif CJK SC}}
}{
  \setCJKmainfont[BoldFont={Heiti SC}]{STSong}
  \setCJKsansfont{Heiti SC}
  \setCJKmonofont{Heiti SC}
  \setmainfont{Times New Roman}
  \newcommand{\kaishu}{\CJKfontspec{STKaiti}}
  \newcommand{\heiti}{\CJKfontspec{Heiti SC}}
  \newcommand{\songti}{\CJKfontspec{STSong}}
}

% ============ 颜色（IB Navy/Grey palette） ============
\usepackage{xcolor}
\definecolor{navy}{HTML}{003366}
\definecolor{deepnavy}{HTML}{0B1F33}
\definecolor{steelgrey}{HTML}{4A5568}
\definecolor{mutedgrey}{HTML}{718096}
\definecolor{lightgrey}{HTML}{F7FAFC}
\definecolor{panelgrey}{HTML}{EDF2F7}
\definecolor{riskred}{HTML}{C53030}
\definecolor{riskamber}{HTML}{D69E2E}
\definecolor{riskgreen}{HTML}{38A169}
\definecolor{riskblue}{HTML}{2B6CB0}
\definecolor{accentblue}{HTML}{2B6CB0}
\definecolor{nvgreen}{HTML}{76B900}
\definecolor{rulegrey}{HTML}{CBD5E0}

% ============ 表格 ============
\usepackage{booktabs}
\usepackage{longtable}
\usepackage{tabularx}
\usepackage{multirow}
\usepackage{array}
\usepackage{makecell}
\usepackage{colortbl}
\renewcommand{\arraystretch}{1.12}
\newcolumntype{Y}{>{\centering\arraybackslash}X}
\newcolumntype{L}[1]{>{\raggedright\arraybackslash}p{#1}}
\newcolumntype{C}[1]{>{\centering\arraybackslash}p{#1}}
\newcolumntype{R}[1]{>{\raggedleft\arraybackslash}p{#1}}

% ============ 图形 ============
\usepackage{graphicx}
\usepackage{tikz}
\usetikzlibrary{shapes,arrows.meta,positioning,calc,backgrounds,fit,decorations.pathreplacing}
\usepackage{pgfplots}
\usepgfplotslibrary{polar}
\pgfplotsset{compat=1.18}
\usepackage[edges]{forest}

% ============ 页眉页脚（报告标题和日期需在main.tex中覆盖） ============
\usepackage{fancyhdr}
\pagestyle{fancy}
\fancyhf{}
\fancyhead[L]{\scriptsize\color{steelgrey} \reporttitle}
\fancyhead[R]{\scriptsize\color{steelgrey} \reportdate}
\fancyfoot[C]{\small\color{steelgrey} \thepage}
\fancyfoot[R]{\tiny\color{steelgrey} \reportdisclaimer}
\renewcommand{\headrulewidth}{0.4pt}
\renewcommand{\footrulewidth}{0.2pt}

% ============ 章节格式 ============
\usepackage{titlesec}
\titleformat{\chapter}[block]
  {\normalfont\LARGE\bfseries\color{navy}}
  {\thechapter}{0.9em}{}
\titlespacing*{\chapter}{0pt}{0.4em}{1.0em}
\titleformat{\section}
  {\normalfont\large\bfseries\color{navy}}
  {\thesection}{1em}{}
\titleformat{\subsection}
  {\normalfont\normalsize\bfseries\color{steelgrey}}
  {\thesubsection}{1em}{}

% ============ 超链接 ============
\usepackage{hyperref}
\hypersetup{
  colorlinks=true,
  linkcolor=navy,
  urlcolor=navy,
  citecolor=navy,
}

% ============ 其他工具 ============
\usepackage{enumitem}
\setlist{nosep, leftmargin=2em}
\IfFileExists{fontawesome5.sty}{\usepackage{fontawesome5}}{}
% ============ Class B 全局兜底（xeCJK + 缺字 fallback 阻断）：
% 当 fontawesome5 因 Times New Roman 缺字 / \FA 命名空间失败时，
% 禁止输出 aExclamationTriangle / aCloud 裸文本，一律降级为 amssymb 原生符号。
% 注：amssymb 在 L112 已加载，此处不重复。所有 \def 在本段落。
\def\faExclamationTriangle{\ensuremath{\textcolor{riskred}{\blacktriangle}}\kern-0.25em}
\def\faCloud{\ensuremath{\clubsuit}\kern-0.1em}
\def\faYenSign{\ensuremath{\yen}}
\def\faIndustry{\ensuremath{\S}}
\def\faCogs{\ensuremath{\ast}}
\def\faFlagCheckered{\ensuremath{\blacktriangleright}}
\def\faMemory{\textbf{[RAM]}}
\def\faServer{\textbf{[DC]}}
\def\faGlobe{\ensuremath{\diamond}}
\def\faStar{\ensuremath{\bigstar}}
\def\faMicrochip{\textbf{[$\mu$P]}}
\def\faLayerGroup{\textbf{[LAY]}}
\def\faCubes{\ensuremath{\boxtimes}}
\usepackage{amssymb}
\usepackage{pifont}
\usepackage{tcolorbox}
\tcbuselibrary{skins,breakable}
\usepackage{caption}
\captionsetup{font=small, skip=4pt}
\usepackage{float}
\usepackage{needspace}
\setlength{\textfloatsep}{8pt plus 2pt minus 2pt}
\setlength{\floatsep}{6pt plus 2pt minus 2pt}
\setlength{\intextsep}{6pt plus 2pt minus 2pt}
\setlength{\abovecaptionskip}{4pt}
\setlength{\belowcaptionskip}{2pt}

% ============ 自定义命令 ============
\newcommand{\rmark}{\textcolor{riskred}{\ding{108}}}
\newcommand{\amark}{\textcolor{riskamber}{\ding{108}}}
\newcommand{\gmark}{\textcolor{riskgreen}{\ding{108}}}
\newcommand{\starfive}{$\bigstar\bigstar\bigstar\bigstar\bigstar$}
\newcommand{\starfour}{$\bigstar\bigstar\bigstar\bigstar$}
\newcommand{\starthree}{$\bigstar\bigstar\bigstar$}
% 资料来源注：使用楷体（STKaiti）而非斜体，避免中文无斜体字形导致的字体回退警告
% \mdseries 避免继承周围 \bfseries 导致 STKaiti 粗体缺失警告
% \sloppy + \hbadness 允许长文件名/URL 在必要时断行
% 局部重定义 \textbf → \heiti：STKaiti 无粗体，用黑体替代强调
\newcommand{\sourcenote}[1]{\par{\scriptsize\mdseries\color{mutedgrey}\kaishu\sloppy\hbadness=9999\def\textbf##1{{\heiti ##1}} 资料来源：#1}\vspace{0.2em}}
\newcommand{\qualitynote}[1]{\par{\scriptsize\mdseries\color{mutedgrey}\sloppy\def\textbf##1{{\heiti ##1}} 数据质量：#1}\vspace{0.2em}}
\newcommand{\evidencetag}[1]{%
  \tikz[baseline=(x.base)]\node[
    xshift=0pt,
    inner xsep=3.5pt,
    inner ysep=1.2pt,
    rounded corners=1.5pt,
    fill=panelgrey,
    draw=rulegrey,
    text=steelgrey,
    font=\scriptsize
  ] (x) {#1};%
}
\newcommand{\stance}[2]{\textbf{\textcolor{#1}{#2}}}

% 信息框
\newtcolorbox{keyinsight}[1][]{
  colback=lightgrey, colframe=navy,
  fonttitle=\bfseries, title={#1}, breakable,
  arc=2pt, boxrule=0.5pt, left=8pt, right=8pt, top=6pt, bottom=6pt
}
\newtcolorbox{riskbox}[1][]{
  colback=red!5, colframe=riskred,
  fonttitle=\bfseries, title={#1}, breakable,
  arc=2pt, boxrule=0.5pt, left=8pt, right=8pt, top=6pt, bottom=6pt
}
\newtcolorbox{dashboardbox}[1][]{
  enhanced, breakable,
  colback=white, colframe=navy,
  boxrule=0.7pt, arc=2pt,
  fonttitle=\bfseries\color{white},
  coltitle=white, colbacktitle=navy,
  title={#1},
  left=9pt, right=9pt, top=7pt, bottom=7pt
}
\newtcolorbox{houseviewbox}[1][]{
  enhanced, breakable,
  colback=lightgrey, colframe=deepnavy,
  borderline west={3pt}{0pt}{accentblue},
  boxrule=0.4pt, arc=2pt,
  fonttitle=\bfseries, title={#1},
  left=9pt, right=9pt, top=7pt, bottom=7pt
}
\newtcolorbox{sourcequalitybox}[1][]{
  enhanced, breakable,
  colback=yellow!6, colframe=riskamber,
  boxrule=0.45pt, arc=2pt,
  fonttitle=\bfseries, title={#1},
  left=9pt, right=9pt, top=7pt, bottom=7pt
}
% exhibitbox：表/图外框。left=7pt+right=7pt=14pt 内边距，before upper 内的 \linewidth 已扣除该内边距。
% 约定：\exhibitboxwidth = 真实内边距宽（tcolorbox 已扣好 padding 的 \linewidth）。
% 由于 tabularx/longtable 宽度参数的语义是「外部总宽（含两侧 \tabcolsep）」，内部列实际可用 = W - 2\tabcolsep。
% 所以 tabularx 直接传 \exhibitboxwidth（总宽 == 内边距宽），而其他用途（resizebox、tikz、p{...}）用 \exhibitwidth = \exhibitboxwidth - 2\tabcolsep。
\newlength{\exhibitboxwidth}
\newlength{\exhibitwidth}
\newtcolorbox{exhibitbox}[1][]{
  enhanced, breakable,
  colback=white, colframe=rulegrey,
  borderline north={1.2pt}{0pt}{navy},
  boxrule=0.35pt, arc=1pt,
  fonttitle=\bfseries\color{navy}, title={#1},
  left=7pt, right=7pt, top=6pt, bottom=6pt,
  before upper={\global\exhibitboxwidth=\linewidth\relax
    \global\exhibitwidth=\dimexpr\linewidth-2\tabcolsep\relax},
}
\newtcolorbox{disclosurebox}[1][]{
  enhanced, breakable,
  colback=lightgrey, colframe=steelgrey,
  boxrule=0.35pt, arc=1pt,
  fonttitle=\bfseries, title={#1},
  left=8pt, right=8pt, top=6pt, bottom=6pt
}

% ============ First-page dashboard cover ============
\newcommand{\astockcover}{%
\begin{titlepage}
\thispagestyle{empty}
\begin{tikzpicture}[remember picture, overlay]
  \fill[deepnavy] (current page.north west) rectangle ([yshift=-2.55cm]current page.north east);
  \fill[accentblue] ([yshift=-2.55cm]current page.north west) rectangle ([yshift=-2.66cm]current page.north east);
  \fill[deepnavy] (current page.south west) rectangle ([yshift=1.2cm]current page.south east);
\end{tikzpicture}
\vspace*{0.62cm}
{\color{white}\sffamily\bfseries\small \reportkicker}\\[0.18cm]
{\color{steelgrey}\sffamily\scriptsize \reportscope \quad|\quad \reporttype \quad|\quad \reportdate}\\[1.0cm]
{\fontsize{25}{30}\selectfont\bfseries\color{navy}\reporttitle}\\[0.25cm]
{\fontsize{13}{17}\selectfont\color{steelgrey}\reportsubtitle}\\[0.75cm]
\begin{dashboardbox}[决策摘要]
\begin{tabularx}{\textwidth}{L{3.0cm}X L{3.0cm}X}
\textbf{报告类型} & \reporttype & \textbf{数据截至} & \reportdatacutoff \\
\textbf{覆盖范围} & \reportscope & \textbf{作者} & \reportauthor \\
\end{tabularx}
\vspace{0.45em}
\begin{houseviewbox}[本机构观点]
\reporthouseview
\end{houseviewbox}
\begin{sourcequalitybox}[证据质量]
\reportquality
\end{sourcequalitybox}
\end{dashboardbox}
\vfill
{\scriptsize\color{steelgrey}\reportdisclaimer}
\end{titlepage}
}

\usepackage{amsmath}
\newcolumntype{P}{>{\raggedright\arraybackslash}X}
\hypersetup{pdfauthor={\reportauthor},pdftitle={\reporttitle}}
\fancyfoot[R]{}

\begin{document}
\astockcover

\begin{abstract}
本报告经最新授权修正后，只保留一个可审计的策略研究目标：在完整、正式指定组合上，争取滚动五年、成本后年化时间加权收益率（TWRR）达到5\%。这不是年度下限、收益保证、概率判断或实施授权。现金分配是总回报和产品尽调的一项事实输入，而非独立的收款目标、覆盖门槛或年度验收项；一笔已结算分配在总回报中只作一次处理。

本次修正取消把首轮产品尽调范围当作永久投资限制的做法。研究宇宙向合规可及的直接股票、ETF、公募REIT、债券与现金工具开放；现有十只产品只是首轮证据队列。任何候选仍须先进入有限、带日期的决策宇宙，完成与其结构相适应的权益研究级尽调，再接受独立估值、风险、流动性、税务/运营和组合层复核。没有资产类别配额必须填满，现金也可以是有效选择。

当前批准可投池与批准组合均为空。正式回放、候选标的风险预算、适当性、账户税费和实时资格未完成前，所有市场交易资产$q=0$；这表示模型未授予市场风险配置权限，并不推定任何账户的实际余额或现金用途。本报告不提供实时价格、买卖数量、交易时点、税务意见、适当性认证、订单或任何收益保证。
\end{abstract}

{\small
\setcounter{tocdepth}{0}
\tableofcontents
\par}
\clearpage

\chapter{投资结论与目标定义}
\begin{houseviewbox}[投委会结论：开放研究宇宙；当前没有获批组合]
经授权修订后，本案只保留滚动五年、成本后年化TWRR达到5\%这一项条件性、事后研究目标。市场策略不能被诚实地包装为年度最低收益或收益保证；目标也不等同于概率判断或实施授权。已结算现金分配只作为总回报核算和产品尽调中的事实输入：它在总回报中只处理一次，其后留存、取出或日后另行考虑配置，均不创造第二层回报。本框架采用比例口径、与资金规模无关，不预设任何账户层事实。

原有的 ETF／REIT 产品菜单和10\%／15\%／40\%／25\%／10\%向量，只能代表首轮研究设计，不能成为永久投资禁区或隐性配置指令。它会在研究开始前排除可能具有更好风险调整回报、现金质量、估值或分散化贡献的直接股票及其他合规工具。自本版起，研究宇宙对所有合规可及的直接股票、ETF、公募REIT、债券和现金工具开放；现有十只产品仅是已完成首轮逐项尽调的队列。

截至本报告日，T3批准可投池和T4批准组合均为空：没有指定直接股票，十只产品也均未完成实时估值、流动性、账户、组合回放和风险预算放行。因此当前动作仍是\textbf{不配置任何市场交易资产}。这表示模型未授予市场风险配置权限；合格受保障流动性仅是未通过资格时的条件性回退类别，不满足收益目标。
\end{houseviewbox}

\begin{exhibitbox}[研究宇宙、可投池与组合不是同一个集合]
\small
\begin{tabularx}{\exhibitboxwidth}{L{2.2cm}L{3.0cm}P P}
\toprule
\textbf{层级} & \textbf{定义} & \textbf{当前状态} & \textbf{能否配置} \\
\midrule
T0 研究边界 & 法律、账户、税务、结算、流动性与授权允许研究的产品格式；不是代码白名单。 & 已定义。杠杆、衍生品、私募或显著低流动性工具需另行书面授权。 & 否。\\
T1 开放研究宇宙 & 可提名直接股票、ETF、公募REIT、主权债与现金工具。现有十只不是排他菜单。 & 开放；所有新提名默认$q=0$。 & 否。\\
T2 决策宇宙 & 针对某一次评审冻结的有限候选，含代码、研究问题、责任人、时间戳和证据路径。 & 十只产品构成首轮完成队列；未冻结直接股票名单。 & 否。\\
T3 批准可投池 & 已完成单项R0--R4、独立证据/估值/风险/运营复核的具体标的。 & 空。 & 仅可进入组合评审。\\
T4 批准组合 & 从T3按边际总回报、估值、风险、重叠、容量和流动性动态选择。 & 空；$q_{portfolio}=0$。 & 只有取得组合审批后才可讨论。\\
\bottomrule
\end{tabularx}
\end{exhibitbox}

\section{目标、证据与当前动作}

\begin{exhibitbox}[目标不是配额，研究通过也不是配置授权]
\small
\begin{tabularx}{\exhibitboxwidth}{L{3.0cm}P P}
\toprule
\textbf{约束／目标} & \textbf{当前证据结论} & \textbf{当前动作} \\
\midrule
每年最低5\%收益 & 未被证明；市场策略不能作最低收益承诺。研究口径为滚动五年、成本后TWRR目标。 & 不把“目标”写成承诺；正式回放完成前不批准实施。 \\
现金分配与税费处理 & 历史分派只能说明产品机制或历史事实，不能替代未来总回报、估值或实施资格。 & 将分配、税费、成本和返还性成分用于总回报归因与产品尽调；不设独立收款门槛。 \\
标的选择自由度 & 直接股票原被预先排除，限制了跨资产比较；开放宇宙只解决提名权，不代表自动投资资格。 & 建立有限、带日期的决策宇宙；每项按其结构完成尽调。 \\
实时实施资格 & 市场、账户、容量、税费、相关性与适当性资料均未完整。 & 每个候选$q=0$；没有下单、份额或建仓时间表。 \\
\bottomrule
\end{tabularx}
\end{exhibitbox}

\section{直接股票不是豁免，也不是永久禁区}

直接股票可以贡献盈利增长、股息增长、估值修复或与既有基金不同的经济暴露；但“高股息”本身不能证明现金可持续，更不能证明其优于 ETF、REIT 或现金工具。每一个拟纳入T3的直接股票必须先在T2名单中按代码登记，并拥有独立、可更新的研究档案：公司及行业证据、审计财务、驱动因素至收入/利润/EPS桥、治理和资本配置、当前价格下的熊/基/牛估值、股息/自由现金流覆盖、催化与证伪、A股自由流通量/停复牌/涨跌停/容量，以及同全部现有候选的发行人集团、行业、因子和共同压力映射。T3的已命名股票数必须与估值及券商材料逐一勾稽；没有匹配材料的股票不得进入T3。

同样的严谨度适用于非股票产品，但模板不能机械照搬。ETF需要基金规则、持仓穿透、跟踪和参考值机制，以及重大证券借贷、衍生品和执行能力；公募REIT需要项目、租户或承购方、运营、债务、资本开支、资产估值和分派桥。产品格式不能替代逐项尽调。

\section{从提名到组合的六道串联门槛}

\begin{enumerate}
  \item \textbf{授权与可及性：}确认产品、账户、税务、结算和流动性属于允许研究范围；
  \item \textbf{研究准入：}冻结有限候选并写明研究问题、数据截止日、责任人和来源计划；
  \item \textbf{单项权益研究级尽调：}完成适配产品结构的R0--R4证据、模型、审稿和渲染记录；
  \item \textbf{独立复核：}证据、估值、风险、流动性、税务和运营控制分别复核；
  \item \textbf{组合审批：}评估边际总回报、估值、分配与现金流质量、相关性、集中度、共同风险、压力损失、容量和适当性；没有强制类别配额；
  \item \textbf{持续监控：}资料过期、重大事件、估值/风险假设失效或流动性恶化，标的立即降为受限或不可增持。
\end{enumerate}

任一门槛缺失，候选保持$q=0$，模型不授予风险配置权限。快速事件研究可以缩短信息收集时间，但不能豁免估值可复算、风险聚合和独立复核。

\section{目前为什么没有批准组合}

本案当前没有一份冻结后的跨资产组合回放，也没有可比较的同日现金、成本、税费、流动性、集中度和风险预算数据。首轮十只产品的逐项研究已说明产品机制和证据缺口，但不等同于T3批准；直接股票尚未命名，更不能以“开放”一词直接进入组合。下一次策略评审必须从T1中冻结T2候选，再以同一套收益、风险和执行标准与现有产品比较。


\chapter{总回报、外部现金流与计量口径}
\section{现金不是第二项目标：采用可处理外部现金流的业绩口径}

本策略只有5\%滚动五年TWRR这一项研究目标。利息、股息、ETF分红和REIT分配不是额外收益层，也不构成独立的现金收款、收益率或覆盖目标。它们的作用有两项：第一，作为持有期间总回报的组成部分；第二，作为研究产品经营质量、资本回收和税费影响的事实输入。无论现金随后留在账户、被取出或未作处理，同一笔现金都不得被重复计入业绩。

策略本身采用\textbf{时间加权收益率（TWRR）}。每次外部现金流都切断一个子期间；子期间起点取前一笔外部现金流后的市值，终点取下一笔外部现金流前市值：

\[
r_j=\frac{V^{pre-flow}_{j,end}}{V^{post-flow}_{j,start}}-1,\qquad
TWRR=\prod_{j=1}^{m}(1+r_j)-1.
\]

已结算分配在发生外部流出前属于\(V\)，不另加一笔“分红回报”；若投资者把该笔资金取出，取出时作为外部现金流切分子期间。投资者自身体验另报\textbf{资金加权收益率（MWRR/XIRR）}：使用每笔真实入金、取现及期末价值的日期和符号求解。TWRR评估策略，MWRR评估资金进出时点；二者不能互换。

\begin{exhibitbox}[总回报台账：分配、处置与外部现金流的处理]
\small
\begin{tabularx}{\exhibitboxwidth}{L{3.4cm}P P}
\toprule
\textbf{项目} & \textbf{业绩记录} & \textbf{控制边界} \\
\midrule
已结算利息、股息、ETF分红、REIT现金分配 & 记录到账日、税前金额、税费、净额、结算凭证及成分标记；在组合价值中反映一次。 & 既不形成独立业绩目标，也不得在TWRR之外再次加回。 \\
卖出ETF或REIT取得的资金 & 作为资产处置、再平衡或外部取现记录。 & 不是新增投资收益；不得被伪装为分配或第二层回报。 \\
净值上涨、浮盈、应计未到账收益 & 随适用估值口径反映于市值，未结算项目单列。 & 不得与已结算分配重复归因。 \\
外部入金、账户划转及取现 & 作为外部现金流切分TWRR子期间。 & 改变投资者的资金体验，不倒推或改写既有策略回报。 \\
\bottomrule
\end{tabularx}
\end{exhibitbox}

返还性成分的台账处理必须在冻结T2组合前声明：在独立确认其经济性质前，应单列、透明记录，且不归类为投资收益；无论最终会计处理如何，都不得重复计入总回报。当前并未冻结任何组合或账户台账，因此不对实际分配作分类判断。

基金分红通常伴随基金资产和净值除权调整。多份REIT公告也说明分派不会降低投资风险或创造额外回报；现金只是财富形态和回报归因的一部分。\sourcenote{\texttt{508055}、\texttt{508068}、\texttt{508077}、\texttt{508080}、\texttt{508096}收益分配公告（R05--R09）；\texttt{510880}年报（R11）。}

\section{分配、税费与成本：用于归因和尽调，不是收益率门槛}

产品研究仍须核对已公告分配、可能的削减、持有人税费、到账费用、基金／券商／托管／再平衡成本以及公司行动。这些输入用于还原历史总回报、识别现金流质量和判断某项资产的经济属性；它们不构成任何最低分配、净现金率、覆盖率或达标公式。历史资料不足、税费未知或成分不清时，正确结论是不能完成该候选的总回报归因或实施尽调，而不是以高收益标签、卖出所得或风险上移补足某个现金缺口。

正式回放只设B0受保障流动性机会成本和B1冻结的T2跨资产组合两个比较基准，并按相同日期、总回报、成本和外部现金流口径执行；在候选、账户工具和成本口径冻结前，任何比较都不具有效力。


\chapter{开放研究宇宙与首轮队列}
\section{首轮产品队列服务于比较，不构成排他菜单}

现有十只产品被完整列入首轮研究，是因为它们曾是策略研究的起点，并不意味着它们构成长期唯一选择。开放研究宇宙允许直接股票、ETF、公募REIT、债券和现金工具在同一轮决策中竞争；但每一次组合评审都必须先冻结有限候选名单。一个未命名的“类别”不能获得估值、总回报或分散化信用，一个基金标签也不能替代底层经济暴露的研究。

\begin{exhibitbox}[首轮已尽调产品队列：用于后续跨资产比较，当前均为$q=0$]
\footnotesize
\begin{tabularx}{\exhibitboxwidth}{L{1.35cm}L{2.45cm}L{2.3cm}P}
\toprule
\textbf{代码} & \textbf{产品} & \textbf{已确认的研究角色} & \textbf{进入T3前仍须完成的事项} \\
\midrule
\texttt{511010} & 五年期国债ETF & 利率曲线暴露，不是受保障流动性。 & 同日价格/参考值、久期/DV01、价差、深度、容量、账户及组合压力。 \\
\texttt{510880} & 上证红利ETF & 红利权益暴露；采矿和金融权重会形成周期与金融因子集中。 & 基金当前成份、TTM实收现金、税后现金桥、发行人集团/因子/REIT交叉风险与容量。 \\
\texttt{512890} & 红利低波ETF & 低波名称不等于低组合风险；金融权重高且近年基金层分配有限。 & 基金当前成份、分配政策、现金质量、与其他权益候选的重叠、估值与容量。 \\
\texttt{508069} & 南京交通REIT & 单项目收费公路现金流。 & 车流、收费、期限、债务、当前估值、流动性和共同风险。 \\
\texttt{508011} & 物美消费REIT & 社区商业租金与租户现金流。 & 租户信用、租约、债务、资本开支、NAV日期与可执行容量。 \\
\texttt{508068} & 北京保障房REIT & 政策租金、出租和收缴。 & 扩募后可比估值、运营、债务、税费和组合层共同风险。 \\
\texttt{508077} & 有巢租赁住房REIT & 租赁住房现金流。 & 同簇住房REIT的相关性、经营KPI、估值、债务和流动性。 \\
\texttt{508055} & 上海租赁住房REIT & 租赁住房现金流。 & 地区/管理人/租赁因子重叠、运营、债务、估值和容量。 \\
\texttt{508080} & 亦庄产业园REIT & 产业园租户现金流；历史价格/NAV日期错配诊断为负。 & 租户集中、出租率、续约、债务、现金流、当前估值和流动性。 \\
\texttt{508096} & 中航京能国际能源REIT & 水光能源现金流。 & NAV口径勾稽、发电/电价、应收与保理、债务、估值和流动性。 \\
\bottomrule
\end{tabularx}
\end{exhibitbox}

这十项已具备产品级资料卡，不等于它们优于直接股票或其他合规工具。对直接股票，后续T2名单中的每一家公司也必须按\texttt{equity-research}标准完成独立深度报告；对ETF和REIT，则用结构适配的产品尽调替代公司模板，而不是降低证据标准。

\section{公募REIT：具有分配机制的研究候选，而非回报承诺}

基础设施REIT在符合分配条件时，规则要求每年将合并后年度可供分配金额的90\%以上以现金形式分配；可供分配金额须在净利润基础上调整，并综合考虑持续发展、偿债能力和经营现金流。该规则说明产品机制，不规定买入成本上的分配率，也不保证二级市场总回报。\sourcenote{公开募集基础设施证券投资基金指引（试行）（R01）。}

\begin{exhibitbox}[已披露现金事实不等于可配置结论]
\footnotesize
\begin{tabularx}{\exhibitboxwidth}{L{1.4cm}L{2.6cm}L{2.2cm}P}
\toprule
\textbf{代码} & \textbf{候选} & \textbf{已披露现金事实} & \textbf{为何仍不能以此批准配置} \\
\midrule
\texttt{508069} & 华夏南京交通高速公路REIT & 2025年化现金流分派率9.19\%；2026年首期每10份CNY1.4980。同期年报披露内部收益率1.44\%。 & 当前价格、TTM分配、收费量、剩余期限、债务、估值和流动性缺失。 \\
\texttt{508011} & 嘉实物美消费REIT & 2025-12-31基准日可供分配CNY35.8781m，拟分配CNY35.2400m。 & 当前市价/NAV、消费物业运营、租户、债务与成交容量缺失。 \\
\texttt{508068} & 华夏北京保障房REIT & 2025年合并可供分配CNY74.8239m；本次每10份CNY0.6000。 & 当前市价/NAV、出租率、租金、资本开支和集中度缺失。 \\
\texttt{508077} & 华夏华润有巢租赁住房REIT & 2025年可供分配CNY50.5599m；本次每10份CNY0.1462。 & 与其他租赁住房候选的行业重叠、当前估值和流动性缺失。 \\
\texttt{508055} & 汇添富上海地产租赁住房REIT & 2026年首期每10份CNY0.5154。 & 当前市价/NAV、运营、债务与跨区域租赁住房相关性缺失。 \\
\texttt{508080} & 中金亦庄产业园REIT & 2026年首期每10份CNY0.9000。 & 当前价格、园区出租率、租户、项目现金流和流动性缺失。 \\
\texttt{508096} & 中航京能国际能源REIT & 2025年度第二次每10份CNY3.3178。 & 发电量、电价、项目现金流、债务、价格/NAV和流动性缺失。 \\
\bottomrule
\end{tabularx}
\end{exhibitbox}

\section{权益及基金候选：ETF与直接股票采用同一决策纪律}

\texttt{510880}的历史分红只能说明ETF可能分配现金，不等同于未来到账：2026年公告每10份CNY1.4300、基准日净值CNY3.1813；2025年报净值增长-0.45\%、五年45.43\%，并披露每10份分配CNY1.4200。\texttt{512890}则在2024年净值增长22.64\%的同时未作利润分配，说明合同条件、基金净值和指数股息率均不能换算为持有人下一年现金。\sourcenote{\texttt{510880} 2026年分红公告及2025年报（R10--R11）；\texttt{512890} 2024年报（R12）。}

因此不应只买ETF或只买股票。直接股票与ETF/REIT只能在同一份跨资产比较包中，按成本税费后的总回报、估值、分配与现金流质量、风险、重叠、容量和证据质量竞争；资产类别、历史分派率与预设权重均不构成胜出理由。产品的分配机制也不规定策略的现金用途：现金到账后是否保留、取出或另行考虑配置，由投资者独立决定，不改变既有总回报的计量。


\chapter{现金流记录与投资者自主用途}
\section{现金流记录与用途相互独立}

本案当前没有获批市场组合，因此没有由本框架指定的账户现金动作。本章只定义已结算分配和外部现金流的记录、归因与边界，绝不构成买入、留存、取出或再配置指令。现金到账不因付款方、除息日、历史分配率或回报目标而触发自动动作；现金用途由投资者独立决定。若日后在独立授权下考虑新的配置，必须作为带日期的新组合审议并通过全部有效资格闸门。

\begin{exhibitbox}[已结算现金：记录、归因与用途边界]
\small
\begin{tabularx}{\exhibitboxwidth}{L{2.2cm}P P}
\toprule
\textbf{步骤} & \textbf{研究动作} & \textbf{控制点} \\
\midrule
1. 入账与对账 & 记录到账日、来源、税前金额、税额、费用、净额、结算凭证及返还性成分标记。 & 公告、应收款和卖出所得不能提前或错误归类为分配；返还性成分不得被称为收益。 \\
2. 总回报归因 & 以统一台账将已结算分配、适用税费与成本纳入TWRR核算一次。 & 留存、取出或后续再配置均不增加第二笔回报，也不改变既有归因。 \\
3. 用途自主 & 投资者可按自身账户、税务、流动性与授权条件决定保留、取出或不作处理。 & 本框架不设默认用途，不把到账现金视为自动配置权限。 \\
4. 另行审议 & 如投资者独立考虑新的配置，才可提交T3候选的边际价值比较。 & $q_{portfolio}$、风险、估值、税费、容量、集中度或相关性任一未通过即不得形成市场配置。 \\
\bottomrule
\end{tabularx}
\end{exhibitbox}

\section{总回报验收、归因与失败状态}

在具有完整实际账户记录、冻结组合及冻结比较基准后，才可做策略业绩验收。总回报采用TWRR；投资者体验另报MWRR/XIRR。正式回放和后续业绩报告均应展示年度与滚动五年TWRR、成本、税费、回撤、波动、下行偏差、换手、未成交、容量、相关性和集中度，并把分配和公司行动纳入同一总回报口径。首次完整五年观察期前，任何“达到／未达到5\%”的判断都应标注为\textbf{尚无完整观察期}；候选基金或单项目历史不能替代策略结果。

\begin{exhibitbox}[实施后才会使用的总回报状态]
\small
\begin{tabularx}{\exhibitboxwidth}{L{3.0cm}P P}
\toprule
\textbf{结果} & \textbf{报告标签} & \textbf{治理结果} \\
\midrule
完整数据下滚动五年TWRR达到或高于5\% & 条件性研究目标达到 & 仍须重做风险、基准、容量和资格审查；不自动扩张风险。 \\
完整数据下滚动五年TWRR低于5\% & 条件性研究目标未达到 & 对账市场／NAV、公司行动、成本、税费、运营与风险暴露；不得通过提高收益风险来“修复”结果。 \\
观察期不足或关键输入缺失 & 尚无完整观察期／不可判定 & 补齐点时序列、成本税费、公司行动和可复现规则；不以单基金历史代替策略回放。 \\
\bottomrule
\end{tabularx}
\end{exhibitbox}

若总回报证据、风险预算、集中度、适当性或实施资格任一项失效，冻结任何增加市场风险的提案，直至相关证据重新通过。这个治理动作不承诺通过交易“修复”结果，也不对现金分配设置独立的年度成功或失败标准。


\chapter{风险预算与目标失败处理}
\section{从固定配比转向风险预算优先}

此前10\%／15\%／40\%／25\%／10\%参考向量包含40\%红利权益ETF和25\%公募REIT，已知风险资产层的规划压力小计为-25.00\%，高于-18\%风险行动线。该算术说明预设类别权重没有通过风险预算；它不应被修补为另一套未经验证的固定配比，更不能据此永久排除直接股票。该向量现仅保留为“被否决的历史研究设计”，用于说明为什么未来组合必须从获批标的的边际风险贡献出发。

\begin{exhibitbox}[历史参考向量的共同压力：用于否决固定配比，不是新组合的损失预测]
\small
\begin{tabularx}{\exhibitboxwidth}{L{3.1cm}P P P}
\toprule
\textbf{历史研究层} & \textbf{历史参考权重} & \textbf{规划压力输入} & \textbf{已知加权贡献} \\
\midrule
受保障流动性 & 10\% & 0\% & 0.00\% \\
五年期国债ETF & 15\% & -5\% & -0.75\% \\
红利权益ETF & 40\% & -45\% & -18.00\% \\
公募REITs & 25\% & -25\% & -6.25\% \\
\midrule
\textbf{已知资产层小计} & \textbf{90\%} &  & \textbf{-25.00\%} \\
未定义防御分散器 & 10\% & 无获批输入 & 不得纳入通过计算 \\
\bottomrule
\end{tabularx}
\end{exhibitbox}

该表不是历史回撤、VaR、预期损失、概率或最大损失。它只说明旧的类别向量不具备激活资格。下一次组合评审需由T3批准可投池的具体标的生成新的压力、相关性、流动性和总回报归因模型；若风险预算不批准，正确结果仍是$q_{portfolio}=0$，而不是把资金转向任何“更高收益”替代品。

\section{开放研究宇宙下的硬性风险控制}

\begin{exhibitbox}[所有资产格式共用的组合约束]
\small
\begin{tabularx}{\exhibitboxwidth}{L{3.15cm}P P}
\toprule
\textbf{控制} & \textbf{适用逻辑} & \textbf{失败结果} \\
\midrule
单项研究完整性 & 每项必须有结构适配的R0--R4档案、当前估值/参考值、现金机制、风险、催化和证伪。 & 不能进入T3。 \\
共同风险映射 & 直接股票、ETF与REIT按发行人集团、行业、因子、利率、客户/租户/承购方、融资和地域/监管共同驱动聚合。 & 未能映射的候选$q=0$。 \\
边际风险贡献 & 新候选须改善或至少符合获批压力损失、集中度、流动性和容量预算；单独高股息不是充分条件。 & 不进入T4。 \\
分配与现金流质量 & 核对分配削减、股利／经营现金流覆盖、返还性成分、税费和成本，以判断经济质量与总回报归因。 & 不得把资产处置所得、杠杆、低信用或强制高收益配置伪装为回报改善。 \\
流动性与实施 & 审核A股自由流通量、涨跌停、停复牌、价差、深度、参与率、最小单位、税费与结算。 & 无可执行性即$q=0$。 \\
资产类别特别授权 & 杠杆、反向/杠杆产品、期货/期权、私募或显著低流动性资产需要独立书面授权。 & 不得由开放宇宙自动纳入。 \\
\bottomrule
\end{tabularx}
\end{exhibitbox}

\section{观察回撤的强制风险动作}

观察回撤仅用带日期、外部现金流调整后的账户市值和高水位计算：
\[
DD_t=\frac{MV_t^{\mathrm{flow\ adjusted}}}{HWM_t^{\mathrm{flow\ adjusted}}}-1.
\]

\begin{exhibitbox}[风险行动线不是被动“复核”]
\small
\begin{tabularx}{\exhibitboxwidth}{L{3.2cm}P}
\toprule
\textbf{触发} & \textbf{强制治理结果} \\
\midrule
$DD_t>-12\%$且其他闸门均通过 & 仅可讨论已在T3且有效$q=1$的候选；不意味着自动下单。 \\
$DD_t\leq-12\%$ & \textbf{谨慎锁定}：所有风险新增暂停；若投资者另行考虑配置，须先重新承保。 \\
$DD_t\leq-18\%$ & \textbf{风险停止}：所有新的市场风险配置均禁止；取消未执行模型批次，形成事件与去风险方案后方可讨论再批准。 \\
达到获批压力容量的90\%（若未来有例外） & 风险停止继续生效，重做持仓、流动性和资金用途风险预算；无机械重入。 \\
\bottomrule
\end{tabularx}
\end{exhibitbox}

这些是新资金的治理权限，不是自动卖出指令，也不保证亏损在阈值停止。悬停、跌停、跳空和低流动性情形仍可能使实际损失超过行动线。

\section{产品结构专项控制不等于类别配额}

公募REIT在进入T4时，仍须满足至少五个独立项目、三个以上管理人/发起人集团和三个以上运营现金流行业；单项目、单一管理人/发起人和共同风险簇必须受组合上限约束。红利ETF与直接股票则需要发行人集团、行业和因子集中度映射。直接股票不因个股数量不同而自动分散，三只租赁住房REIT也不因代码不同而天然独立。

未来可由风险预算为单一发行人、单一基金、单一REIT项目、行业/因子/现金流风险簇分别设定上限，但本案没有获批的候选组合或损失容量，因此不虚构数值上限。唯一有效的当前控制是：全部市场候选保持$q=0$，直至独立风险复核完成。


\chapter{资格包与条件性实施}
\section{实施资格是乘法闸门，不是交易清单}

本报告不提供下单价格、数量、建仓日期或分批计划。未来每一项具体候选仅在以下乘法闸门全部为1时才达到“可供独立组合审议”的状态：
\[
q_{effective,i}=q_{portfolio}\times q_{risk\ state}\times q_{research,i}
\times q_{valuation,i}\times q_{market,i}\times q_{concentration,i}
\times q_{account,i}.
\]
任一项为0，结果即为0。目前T3为空、$q_{portfolio}=0$，因此所有市场候选$q=0$。开放研究资格、历史分派、产品标签、单个公告或部分同日行情均不推定其他项通过。

\begin{exhibitbox}[不同结构的候选必须补齐的当前数据包]
\small
\begin{tabularx}{\exhibitboxwidth}{L{2.65cm}P P}
\toprule
\textbf{候选结构} & \textbf{单项研究与当前实施证据} & \textbf{失败路径} \\
\midrule
直接A股股票 & 公司/行业和治理档案、审计财务、收入至EPS桥、当前价格熊/基/牛估值、股利/自由现金流覆盖、催化与证伪、自由流通量、停复牌/涨跌停、价差、深度和容量。 & 未完整通过时保持在T1/T2；全部单项及组合闸门通过后才可进入T3/T4审议。 \\
ETF/基金 & 当前文件、指数/授权规则、持仓与重叠、当前价格/IOPV或参考值、跟踪和分配、重大证券借贷/衍生品、价差、深度、容量、最小单位与账户。 & 候选$q=0$；不得因标签或指数股息率自动替换。 \\
公募REIT & 项目、租户/承购方、运营KPI、TTM已结算现金、NAV日期与项目事件、债务/资本开支、税费、相关性、价差、深度和容量。 & 候选$q=0$；不得以高历史分派率强行填满额度。 \\
主权债/现金工具 & 发行人或保护机制、期限、到期/提前支取条款、现金流、税费、账户及集中度。 & 不作为已满足5\%总回报目标的假设。 \\
\bottomrule
\end{tabularx}
\end{exhibitbox}

\section{动态市场质量和容量约束}

交易价格相对适用参考值的偏离必须低于候选估值容忍度与20日历史保护线中的更严格者；价差必须低于全成本预算与20日相对价差保护线中的更严格者。拟议金额不得超过目标缺口、当日成交额、20日中位成交额和可执行深度所导出的最严格容量上限。每个输入必须同币种、为正且可追溯。REIT披露NAV不是实时IOPV，直接股票也不能以一个收盘价替代完整的停牌、涨跌停和容量评估。

缺失、过期、停牌、异常估值偏离、无可执行深度、账户税费未对账或任何集中度超限均令相应$q=0$。未通过的候选不获得市场配置权限；它可以在下一轮开放研究中与新的合规候选比较，但不能因为历史分配、回报目标或产品标签自动转入任何高收益、低信用、相关ETF、直接股票或“补仓”批次。

\section{从T3到T4的独立审批记录}

每一项拟进入组合的候选须保留一份带日期的记录：来源定位与时间戳、原始记录哈希、研究档案版本、估值和熊/基/牛假设、税后费后现金输入、流动性/容量计算、共同风险映射、拟议边际权重、组合压力变化、有效$q$、独立数据/估值/风险/运营复核身份以及失败时的现金回退位置。直接股票还须在T2/T3清单中列明代码，并与当前估值行和券商／一致预期材料逐一匹配；声明的T3股票数量必须等于该匹配清单数量。

审批顺序是：数据核验、单项研究复核、独立风险/集中度复核、投资组合治理与适当性复核，再到组合所有者确认是否进入审议。最后一步不授权订单、不指定经纪商、不设定实时价格，也不创建自动交易。

\section{现金流记录与投资者自主用途}

\begin{enumerate}
  \item 记录每笔\textbf{已结算}现金分配的日期、来源、总额、税、费用、净额、结算证据和返还性成分标记；
  \item 按预先冻结的台账规则单列返还性成分；无论处理方式如何，均不得称为收益或重复计入总回报；
  \item 收款后由投资者按自身账户、税务、流动性和授权条件决定保留、取出或不作处理；本报告不规定默认用途；
  \item 只有在投资者独立考虑新的配置时，才可在定期评审中比较有效$q=1$的T3候选；收款本身不触发自动配置；
  \item 卖出资产所得属于处置或再配置，不得伪装为分配、独立回报或目标达成。
\end{enumerate}

\section{今天与未来的边界}

今天没有市场交易动作。完成资料包只会形成一份可由独立风险与投资组合负责人审议的资格记录，并不自动授权交易。若之后发生投资者自主的实际交易决定，仍须在当时用真实账户、市场状态、流动性、停复牌、涨跌停和税费条件作独立判断。


\chapter{回报证据与实施批准门}
\section{已有的历史事实，不是策略业绩}

现有一级资料只支持若干单基金／单项目事实：\texttt{511010} 2025年报披露2023、2024、2025年度净值增长分别为3.28\%、5.95\%、0.59\%，三年和五年累计净值增长分别为10.07\%和16.70\%；\texttt{510880}同期披露年度净值增长7.41\%、19.65\%、-0.45\%，三年和五年累计为27.93\%和45.43\%。\texttt{508069}披露9.19\%年化现金流分派率和1.44\%内部收益率。它们说明分派、净值回报和内部收益率不是同一指标。\sourcenote{\texttt{511010} 2025年报（R14）；\texttt{510880} 2025年报（R11）；\texttt{508069} 2025年报（R02）。}

\begin{exhibitbox}[历史记录能证明什么，不能证明什么]
\small
\begin{tabularx}{\exhibitboxwidth}{L{3.2cm}P P}
\toprule
\textbf{已归档证据} & \textbf{可支持的结论} & \textbf{明确不能支持的结论} \\
\midrule
\texttt{511010}、\texttt{510880}年度报告 & 个别基金在指定期间的净值和分派历史；个别年份可低于5\%。 & 未来T2跨资产组合的历史回报、预期收益、命中率或年度最低收益。 \\
\texttt{508069}年报及REIT分派公告 & 单项目现金分派率与总回报／IRR可分离。 & 五项目REIT篮子的分派、相关性、估值或承受能力。 \\
REIT规则与产品公告 & 分派机制、产品风险和不保证收益边界。 & 持有人明年到账现金、税后收益、现金用途或可交易价格。 \\
\bottomrule
\end{tabularx}
\end{exhibitbox}

因此，回报证据门状态为\textbf{FAIL／NOT RUN／BLOCKED}。不能把若干不同期间、不同计价口径的基金或个股历史加权拼接后称作“回测”；那会忽略T2最终成分、公司行动、费用、税费、滑点、容量、停牌与涨跌停，也会掩盖未冻结的候选与替换规则。

\section{正式回放与基准必须如何完成}

在实施审议前，必须预先冻结：本轮T2直接股票、ETF、REIT、债券与现金工具候选及其替换史、点时权重、总回报与分派／公司行动数据、成本税费和成交规则。回放使用真实A股交易约束，包括最小单位、T+1、停牌、涨跌停无成交、价差、参与率和规则化再平衡；并报告年度与滚动五年TWRR、5\%总回报目标的观察期状态、回撤、波动、下行偏差、换手、成本、未成交、容量、相关性和集中度。分配、税费和公司行动仅在同一总回报口径内归因，不另设独立分配考核。

同时冻结B0受保障流动性机会成本和B1本轮T2跨资产组合两个基准，保持相同日期、估值时点、总回报处理和成本口径。有限历史即使表现良好，也不能单独给出未来5\%的概率；若日后建立概率模型，还必须披露样本、制度处理、分布和相关性假设、参数不确定性与样本外验证。

\section{实施批准矩阵}

\begin{exhibitbox}[四项实施门均为阻断状态]
\small
\begin{tabularx}{\exhibitboxwidth}{L{3.1cm}P P}
\toprule
\textbf{实施门} & \textbf{当前状态} & \textbf{解除条件} \\
\midrule
回报证据 & BLOCKED：正式回放未运行。 & 可复现、点时、成本税费后回放和B0/B1比较报告，经独立无重复计数复核。 \\
风险预算 & NOT APPROVED：旧固定向量的已知压力-25.00\%与-18\%行动线不一致。 & 冻结T2候选，重算候选组合压力、确认损失容量和期限，并取得独立批准。 \\
适当性 & PENDING：未确认账户用途、期限、取现约束、税费和既有持仓重叠。 & 账户与用途资料完整，并由投资者自行确认。 \\
执行资格 & NOT AUTHORIZED：所有$q=0$，无同一时段市场和容量资料。 & 产品、市场、集中度、账户和容量乘法闸门全部为1；仍需独立实施决定。 \\
\bottomrule
\end{tabularx}
\end{exhibitbox}

这四项不会因一份分配公告、收益目标或价格折价自动解除。研究框架通过只表示问题、证据边界和拒绝路径已被清楚记录；它不表示可以开始建仓。

\section{本次投委会授权边界}

本次仅允许继续以下研究工作：从开放宇宙冻结T2候选与替换规则、为每项完成结构适配的权益研究级档案、补齐可复现总回报／分配／公司行动数据、建立B0/B1基准、完成候选分配、账户成本与税费对账，并提交独立风险预算与适当性材料。所有工作都应保留时间戳、原始来源和失败记录。

本次\textbf{不批准}任何直接股票、ETF、REIT、债券或分散器的市场建仓，不批准把受保障流动性替换为收益产品，不批准依据历史分派自行推定未来总回报，也不批准把5\%目标作为年度承诺对外表述。任何未来实施提案必须先冻结新的T2候选并重新进入逐项研究和组合实施门，而不是引用本报告替代当时的投资决定。


\appendix
\begingroup
\small
\chapter{名词解释与计量约定}
\section{阅读口径}

本章解释报告中的产品、计量和控制术语，便于阅读，不替代基金合同、招募说明书、交易规则、账户条款或税务意见。除非四项实施闸门均已通过，以下定义均不构成任何产品的买卖、持有或现金用途指令。

\footnotesize
\begin{longtable}{p{2.9cm}p{7.2cm}}
\toprule
\textbf{术语} & \textbf{本报告中的含义与边界} \\
\midrule
\endfirsthead
\toprule
\textbf{术语} & \textbf{本报告中的含义与边界} \\
\midrule
\endhead
\bottomrule
\endfoot

\multicolumn{2}{l}{\textbf{产品与持有结构}} \\
\addlinespace[0.25em]
公募基础设施REITs（REITs） & REIT是\textit{Real Estate Investment Trust}的国际通用简称。本报告中的“公募REITs”特指公开募集基础设施证券投资基金，并不泛指境外REIT、房地产开发商股票或其他高分红基金。按R01所列产品特征，80\%以上基金资产投向基础设施资产支持证券并持有其全部份额，基金再通过特殊目的载体取得基础设施项目的完全所有权或经营权利，采取封闭式运作；投资者持有的是基金份额。该结构不构成现金分配或二级市场回报保证。 \\

基础设施资产支持证券／特殊目的载体 & 公募REIT持有基础设施项目权益的法律与交易载体。本报告使用该术语说明持有结构，不把它等同于投资者对项目现金流、项目价值或退出价格的保证。 \\

封闭式运作 & 本报告中的基础设施基金按封闭式方式运作；符合条件并经交易所审核同意上市后，基金份额可在交易所交易。二级市场价格可能与披露的净资产值不同。具体存续、交易与终止安排以基金合同和公告为准。 \\

ETF & 交易型开放式基金。本报告的\texttt{511010}和\texttt{510880}均为ETF研究锚，属于市场交易资产而非“受保障流动性”。ETF的历史分红、指数股息率或净值表现不自动等于未来实际到账现金，也不等于本策略的总回报。 \\

开放研究宇宙（T1） & 在授权、账户、税务、结算与流动性边界内，可被提名研究的直接股票、ETF、公募REIT、债券与现金工具集合。它不是持仓名单、推荐清单或自动准入名单；每项默认保持\(q=0\)。 \\

决策宇宙（T2） & 为一次具体比较而冻结的有限候选集。每一项必须有代码或身份、研究问题、数据截止日、责任人和证据路径。开放宇宙保证选择自由，T2保证比较可审计。 \\

批准可投池（T3）／批准组合（T4） & T3是完成结构适配的R0--R4研究并经独立证据、估值、风险和运营复核的具体标的；进入T3只代表可参与组合评审。T4是从T3按边际总回报、估值、风险、重叠、容量和流动性形成的获批组合。两者均不代表存在必须填满的资产类别配额。 \\

直接股票尽调 & 对单一发行人开展公司、行业、财务、盈利桥、治理、现金流/股利、当前估值、催化/证伪、A股流动性和组合重叠研究。它与ETF/REIT尽调同等严格，但证据模块因产品结构不同而不同。 \\

\addlinespace[0.25em]
\multicolumn{2}{l}{\textbf{现金与业绩计量}} \\
\addlinespace[0.25em]
可供分配金额 & 基础设施REIT判断现金分配基础的金额。报告仅采用R01所列规则口径：该金额在净利润基础上调整，并综合考虑项目公司持续发展、偿债能力和经营现金流。它不是净利润本身、不是持有人实际到账净现金，也不能直接当作按买入成本计算的收益率。 \\

现金分配／已结算现金记录 & 基金向持有人支付的现金。报告记录已结算、可核对的税前金额、税费、到账费用和净额，以支持总回报归因与产品尽调；卖出份额所得、未实现收益、净值上涨、外部入金和未结算应收款不得被误分类为分配。可能含返还性成分的分配须单列标记，并依预先声明的台账口径处理，不得称为收益或在总回报之外再加一次。 \\

TTM & \textit{Trailing Twelve Months}，即滚动过去12个月。它可用于查看产品历史分配、经营或估值数据，但不是下一年度分配或总回报预测，也不是本策略的独立验收目标。 \\

总回报 & 本报告将市价／NAV变动、现金分配、税费、交易成本和再平衡影响放入同一可对账口径。现金分配只在该口径内处理一次，不是额外叠加的一层收益或独立业绩目标。 \\

TWRR & \textit{Time-Weighted Rate of Return}，即时间加权收益率。每次外部入金、取现或划转均切分一个子期间，用于评估策略本身，不让投资者资金进出时点扭曲策略回报。 \\

MWRR／XIRR & 资金加权收益率／按日期的内部收益率。它以实际入金、取现和期末价值及其日期计算，用于反映投资者自身的资金体验；不能替代TWRR评价策略。 \\

返还性成分 & 可能反映资本返还而非经营或投资收益的现金组成部分。报告要求将其单列并在经济性质得到独立确认前排除出投资收益归因；最终处理不得造成总回报重复计数。 \\

NAV／IOPV／折溢价 & NAV是基金最近披露的净资产价值；IOPV是盘中参考估值；折溢价是市场价格相对适用参考价值的偏离。三者均不等于可成交价格。本报告不将REIT最新披露的NAV视为实时IOPV；缺少同一时段的价格、参考值和流动性资料时，资格闸门不能通过。 \\

\addlinespace[0.25em]
\multicolumn{2}{l}{\textbf{实施与风险控制}} \\
\addlinespace[0.25em]
资格闸门\(q\)／组合闸门\(q_{portfolio}\) & 报告用于表达实施资格的控制值。\(q=0\)表示该资产层因资料缺失、过期或未通过条件而不得配置；\(q=1\)只表示该层完成规定核验，仍须服从组合风险预算和独立实施决定。当前所有市场资产及组合均为\(q=0\)。 \\

受保障流动性 & 本报告的回退类别，指同时满足账户、产品、期限和相关保护条件的现金或直接国债类安排。它承接未获资格的金额，但不等于整套组合保本，也不保证达到5\%总回报目标。 \\

当前状态／历史参考向量 & “当前状态”是本报告日的模型动作：所有市场资产为\(q=0\)，T3/T4为空，不推定任何账户层事实。旧的10\%／15\%／40\%／25\%／10\%向量仅保留为被否决的首轮研究设计，用于说明固定类别配比的风险缺陷；它不是交易指令或未来配额。 \\

风险预算／共同压力 & 用具体候选的压力输入、相关性、集中度与流动性进行的治理测试，用于决定是否可进入实施审议；不是历史回撤、VaR、概率预测或最大损失承诺。旧向量的已知资产层共同压力小计为-25.00\%，高于-18\%风险行动线，因此不能作为固定配比启用。 \\

外部现金流调整后回撤\(DD_t\) & 相对外部现金流调整后高水位的账户市值跌幅。-12\%和-18\%是本报告对新增风险资金的治理行动线，不是自动卖出、止损保证或损失上限。 \\

风险簇／独立项目 & 本报告的内部集中度控制。基金代码不同不自动代表独立；需同时检查项目、管理人／发起人、行业、地区、主要租户或承购方、融资和相关性。至少五个独立项目只是最低必要条件，并不等于监管规定的“买满五只即可分散”。 \\

B0／B1 & 正式回放所需的两个预先冻结比较基准：受保障流动性机会成本和T2的点时跨资产组合。当前两者均未形成有效可比结果，不能用其推定未来5\%回报的概率。 \\
\end{longtable}
\normalsize

\sourcenote{公募基础设施REIT的法定结构和分配机制依据证监会《公开募集基础设施证券投资基金指引（试行）》（R01）。其余术语用于说明本报告既有的业绩、现金台账、风险和资格控制口径；不新增任何收益、现金或实施结论。}


\chapter{来源、数据质量与重要披露}
\section{来源、数据质量与边界}

本报告优先使用中国证监会、上海证券交易所和基金管理人披露的原始PDF，并已于2026年8月12日归档。每条主张均配有来源编号、原始文件哈希和适用边界；读者不应将历史分配或公告金额外推为未来收益率。

\begin{longtable}{p{1.1cm}p{3.4cm}p{5.7cm}}
\toprule
\textbf{ID} & \textbf{来源} & \textbf{在本报告中的有限用途} \\
\midrule
R01 & 证监会REIT指引 & 90\%现金分配机制及其非收益率下限属性。 \\
R02--R09 & 上交所／基金管理人REIT年报和分配公告 & 历史或已公告现金事实、候选池与风险边界；不用于前瞻收益率。 \\
R10--R12 & 上交所／基金管理人ETF公告和年报 & 红利ETF的条件性现金机制、历史净值与无分配反例。 \\
R13 & \texttt{511010}产品资料概要 & 五年国债ETF身份与不保证最低收益的边界。 \\
R14 & \texttt{511010} 2025年年报 & 2023--2025净值增长、三年／五年累计表现与分配记录；仅是单基金历史，不是组合回放或预期收益。 \\
\bottomrule
\end{longtable}

\qualitynote{资料质量为Tier-1产品、规则和公告资料。它们支持历史或已公告事实和产品机制，但不支持未来收益率、年度现金下限、策略命中率或概率。当前价格、IOPV/NAV、折溢价、交易量、价差、具体税费和实际到账现金均未取得；正式组合回放、风险预算例外和适当性记录亦未完成。该缺失使全部市场交易资产的执行资格保持$q=0$。}

\begin{disclosurebox}[重要披露]
\footnotesize
本报告为研究与教育用途。它不构成投资建议、适当性意见、税务意见、订单、收益保证或资产价值保证。5\%是条件性的滚动总回报研究目标，而不是每年最低收益承诺。基金分红会影响基金资产或净值，不能被理解为总回报之外的额外收益或独立业绩目标。
\end{disclosurebox}

\endgroup

\end{document}
